\section{Il Processo di Compilazione}
Il codice sorgente viene trasformato in binario attraverso diverse fasi:

\begin{enumerate}
    \item \textbf{Preprocessore:} Gestisce inclusioni e macro. (\texttt{\#define}, \texttt{\#if}, \dots)
    \item \textbf{Compilatore:} Traduce il linguaggio ad alto livello in Assembly.
    \item \textbf{Assembler:} Converte Assembly in codice macchina rilocabile.
    \item \textbf{Linker/Binder:} Unisce file oggetto e librerie.
    \item \textbf{Loader:} Converte indirizzi rilocabili in assoluti e carica in memoria.
\end{enumerate}

\subsection{Analisi del Compilatore}
\begin{itemize}
    \item \textbf{Analisi Lessicale:} Rimozione spazi e raggruppamento token.
    \item \textbf{Analisi Sintattica/Semantica:} Controllo struttura e tipi di dati.
    \item \textbf{Ottimizzazione:} Miglioramento del codice per velocità o spazio.
    \item \textbf{Generazione:} Produzione del codice target.
\end{itemize}

\section{Interfacciamento I/O}
Le periferiche sono mappate in memoria (\textit{Memory Mapped I/O}) tramite registri specifici:
\begin{itemize}
    \item \textbf{Datenregister:} Contiene i dati da elaborare.
    \item \textbf{Steuerregister:} Configurazione della periferica.
    \item \textbf{Statusregister:} Indica lo stato corrente dell'I/O (es. Busy, Ready).
\end{itemize}


